%------------------------------------------------------
% Cover Letter -- NGUYEN KHAC BAO
% Target: AI Engineer Intern (Game Development) @ Gameloft Saigon
% Compile: pdflatex CoverLetter_Nguyen_Khac_Bao_Gameloft.tex
%------------------------------------------------------

\documentclass[11pt, a4paper]{article}

\usepackage[T1]{fontenc}
\usepackage[utf8]{inputenc}
\usepackage[a4paper, top=0.9in, bottom=0.9in, left=1in, right=1in]{geometry}
\usepackage{hyperref}
\usepackage{parskip}
\usepackage{microtype}

\hypersetup{
    colorlinks=true,
    urlcolor=black,
    linkcolor=black,
    pdfborder={0 0 0}
}

\setlength{\parindent}{0pt}
\setlength{\parskip}{8pt}
\pagestyle{empty}

\begin{document}

%---------- SENDER ----------
\begin{flushleft}
\textbf{Nguyen Khac Bao}\\
Ho Chi Minh City, Vietnam\\
\href{mailto:nguyenkhacbaowork564@gmail.com}{nguyenkhacbaowork564@gmail.com} \,\textbullet\,
(+84) 039 586 4429 \,\textbullet\,
\href{https://github.com/NguyenKhacBao564}{github.com/NguyenKhacBao564}
\end{flushleft}

\vspace{4pt}
\hfill May 9, 2026

\vspace{8pt}

%---------- RECIPIENT ----------
\begin{flushleft}
\textbf{Recruitment Team}\\
Gameloft Saigon\\
Ho Chi Minh City, Vietnam
\end{flushleft}

\vspace{4pt}

\textbf{Re: Application for AI Engineer Intern (Game Development)}

\vspace{8pt}

%---------- BODY ----------
Dear Hiring Team,

I am writing to apply for the AI Engineer Intern position at Gameloft Saigon. As a fourth-year Information Technology student at the Posts and Telecommunications Institute of Technology (PTIT) focused on Computer Vision, I am drawn to this role because it sits at the intersection of two things I genuinely care about: building reliable datasets at production scale, and seeing that work power a real product that users---including children using Applaydu---rely on every day.

Over the past year I have completed three personal CV projects, and the common thread among them is dataset preparation rather than model novelty. In my Real-Time Face Mask Compliance Detection System, I converted PWMFD from PASCAL VOC to YOLO format (9{,}205 images, 18{,}528 labelled instances), built audit scripts that surface malformed labels and class-distribution imbalance, and trained YOLOv8n to a validation mAP@50 of 0.970---while transparently documenting the underrepresented \textit{incorrect\_mask} class and its impact on per-class recall. I have also planned a dataset extension with hand-occlusion subclasses, because I noticed that real cameras often capture users mid-gesture and a model trained only on clean frames will fail silently in deployment.

What I would like to mention modestly is a small habit of building lightweight tools when existing options do not quite fit the workflow. For my FallGuard project, I built a Streamlit-based labeling assistant that auto-suggests 13-keypoint pose skeletons via YOLO-Pose, then routes each frame through a human-in-the-loop review UI with confidence display, label correction, uncertain/exclude actions, and manual keypoint editing. It is not a replacement for industry tools, but the experience taught me that small productivity-focused tooling---designed around the actual review bottleneck---can shorten an annotation cycle meaningfully and produce cleaner data than brute-force manual labeling.

I would also like to be transparent about one current limitation: my spoken English is at a basic level. I can read technical documentation and follow written communication comfortably, but my conversational fluency is limited to simple words and sentences. I am actively working to improve it, and I prefer to ask for clarification rather than guess when something is unclear.

I am available to commit full-time (5 days per week) for the duration of the internship and can start at the team's convenience. I would be glad to discuss how I can contribute to dataset quality, annotation workflows, and edge-case handling for the team's ongoing CV projects. Thank you for considering my application.

\vspace{12pt}

Sincerely,

\vspace{18pt}

\textbf{Nguyen Khac Bao}

\end{document}
